% ===================================================================
%  RESUMEN Y SUMMARY
% ===================================================================
\newpage
\thispagestyle{fancy}

\addcontentsline{toc}{chapter}{Resumen}
{\Large\textsc{\textcolor{nar}{Resumen}}}
\vspace{0.2cm}

Producir una retransmisión en directo con varias cámaras, transiciones, rótulos y
grabación simultánea ha requerido históricamente equipamiento de radiodifusión
profesional con costes de adquisición de varios miles de euros y equipos técnicos
especializados desplazados al lugar del evento. Este modelo resulta inaccesible para
comunidades universitarias, organismos públicos y eventos de mediana escala que
necesitan producción de calidad sin presupuesto \textit{broadcast}.

Este Trabajo Fin de Grado presenta el diseño e implementación de Voctomix~2.0, una
plataforma modular de código abierto para la producción remota de vídeo en tiempo
real. Construida sobre GStreamer y Python, sigue una arquitectura cliente-servidor en
la que el núcleo de procesamiento puede ejecutarse en un servidor mientras el
realizador opera la interfaz de control desde cualquier equipo de la red. Las
contribuciones principales del trabajo son: un sistema de composites avanzados
(pantalla completa, imagen dentro de imagen y visión lateral), un módulo de
transiciones animadas entre fuentes, soporte para overlays gráficos dinámicos con
animaciones de entrada y salida, un gestor de estados de emisión con tres modos
(en antena, pausa y sin señal), sincronización automática del audio con la fuente de
vídeo activa, y un sistema de telemetría que registra todos los eventos de la sesión
en tiempo real y los publica en un broker RabbitMQ mediante el protocolo AMQP.

Para facilitar el despliegue reproducible del sistema completo, se ha desarrollado un
entorno contenerizado con Docker~Compose que arranca todos los servicios con un único
comando, y un despliegue alternativo sobre Kubernetes para entornos de mayor escala.

El sistema ha sido validado en dos escenarios reales: dos equipos físicos en red
local y despliegue completo en contenedores. Los resultados confirman una latencia de
conmutación mediana de 1{,}5~ms en Docker y 1{,}8~ms en Kubernetes, 30 veces por
debajo del umbral de la recomendación ITU-R~BT.1359-1, con un uso de CPU inferior al
40\,\% con cuatro cámaras activas. Voctomix~2.0 ha sido integrado en el proyecto
europeo de investigación CyberNEMO de la UPM, donde opera como plataforma de
producción audiovisual en un entorno de validación activo.

\vspace{0.5cm}
\addcontentsline{toc}{chapter}{Palabras clave}
{\Large\textsc{\textcolor{nar}{Palabras clave}}}
\vspace{0.2cm}

producción audiovisual en directo, mezclador de vídeo, software libre, GStreamer,
Docker, Kubernetes, RabbitMQ, telemetría, \textit{streaming}, AMQP.

\newpage
\thispagestyle{fancy}

\addcontentsline{toc}{chapter}{Summary}
{\Large\textsc{\textcolor{nar}{Summary}}}
\vspace{0.2cm}

Producing a live broadcast with multiple cameras, transitions, graphics and
simultaneous recording has historically required professional broadcast equipment
costing several thousand euros and specialised technical crews deployed on site.
This model is out of reach for university communities, public organisations and
medium-scale events that need quality production without a broadcast budget.

This Bachelor's Thesis presents the design and implementation of Voctomix~2.0, a
modular open-source platform for remote video production in real time. Built on
GStreamer and Python, it follows a client-server architecture in which the processing
core can run on a server while the director operates the control interface from any
machine on the network. The main contributions of the project are: an advanced
composite system (fullscreen, picture-in-picture and side-by-side), an animated
transition module between sources, support for dynamic graphical overlays with
blend-in and blend-out animations, an output state manager with three modes (live,
pause and no-signal), automatic audio synchronisation with the active video source,
and a telemetry system that records all session events in real time and publishes them
to a RabbitMQ broker via the AMQP protocol.

To facilitate reproducible deployment of the full system, a containerised environment
using Docker~Compose has been developed that starts all services with a single
command, along with an alternative deployment on Kubernetes for larger-scale
environments.

The system has been validated in two real scenarios: a two-machine deployment over a
local network and a full container-based deployment. Results confirm a median
switching latency of 1.5~ms in Docker and 1.8~ms in Kubernetes, 30 times below the
ITU-R~BT.1359-1 threshold, with CPU usage below 40\,\% with four active cameras.
Voctomix~2.0 has been integrated into the CyberNEMO European research project at
UPM, where it operates as the audiovisual production platform in an active validation
environment.

\vspace{0.5cm}
\addcontentsline{toc}{chapter}{Keywords}
{\Large\textsc{\textcolor{nar}{Keywords}}}
\vspace{0.2cm}

live audiovisual production, video mixer, free software, GStreamer,
Docker, Kubernetes, RabbitMQ, telemetry, streaming, AMQP.

\newpage
